% Generated from the matching Typst scaffold by
% scripts/migrate_typst_report_to_latex.py (prose drafted 2026-06-25).
% NOTE: LightMem (2510.18866) and Mem0 (2504.19413) were [VERIFY+ADD] in the scaffold but
% are NOT in references.bib -- omitted to avoid uncited claims. Add to bib + cite if desired.

% ── 3.1  Memory Management and Forgetting in LLM Agents ──────────────────
% PURPOSE: Survey agent-memory systems thematically; demarcate forgetting layers.
% ─────────────────────────────────────────────────────────────────────────

\section{Memory Management and Forgetting in LLM Agents}

\subsection{Agent memory systems and the store-everything bias}

The first wave of agent-memory systems treats memory as something to grow, not to prune.
MemGPT~\cite{packer2024} borrows the operating-system metaphor, paging information between
a bounded context and external storage so the agent can ``recall'' beyond its window; the
mechanism is elegant, but it pages material in and out rather than deciding what is no
longer worth keeping, so the store only grows. MemoryBank~\cite{zhong2024} does introduce
forgetting --- an Ebbinghaus-style decay over dialogue memories --- but it is built for
conversation and, like most such systems, does not hold retrieval constant across the
variants it compares, so any benefit blends storage with search. A-MEM~\cite{xu2025} adds
structure, organising memories into self-linking notes, yet the design remains
accretion-biased and offers no controlled comparison of \emph{whether} to retain.
HippoRAG~\cite{gutierrez2024hipporag} improves recall with a hippocampally-inspired index,
but indexing is orthogonal to the retention question --- it changes how the store is
searched, not how large it should be allowed to grow. Across this family the default is
accumulation, and storage is never isolated from retrieval.

\subsection{Explicit forgetting and consolidation}

A smaller body of work makes forgetting a first-class operation. The closest to this thesis
is Alqithami's six-policy comparison~\cite{alqithami2025}, which sweeps a comparable breadth
of forgetting behaviours in a privacy-aware cognitive-memory architecture. It is, however,
conversational, and --- decisively for our purposes --- it does not hold retrieval constant
across the compared policies, so its conditions differ in both what they store and how they
search. It demonstrates that forgetting policies are worth comparing without isolating the
forgetting decision itself, which is precisely the control this thesis adds.

\subsection{KV-cache eviction: a different layer}

Finally, a related-sounding literature on cache eviction --- H2O's heavy-hitter
oracle~\cite{zhang2023h2o} and StreamingLLM's attention sinks~\cite{xiao2024streamingllm}
--- ``forgets'' at the level of the attention key--value cache, dropping tokens to keep
decoding cheap over long inputs. It is important to demarcate this clearly: that work
operates inside a single forward pass on the token/attention layer, whereas the six policies
here act on the application-level \texttt{MemoryStore} of past tasks. The two are
complementary rather than competing --- one bounds the context, the other bounds the
experience base --- and conflating them would misattribute footprint savings to the wrong
mechanism.
